% 第5章 5.4.2：仿射核 1+xz 在 [-1,1] 上，系数平面 (a,b)（f=a+bx）中的三个单位球（解析式）。
% 核范数 ||f||_{H_K}^2=a^2+b^2（圆，与测度无关）；
% L^2(dz/2)：a^2+b^2/3=1（半轴 1 与 sqrt3=1.732）；L^2(dz)：2a^2+2b^2/3=1（半轴 0.707 与 1.225）。
% 点 f=1+x：核范数平方 2，L^2(dz/2) 范数平方 4/3。
\begin{tikzpicture}[font=\figurefont\footnotesize,x=17mm,y=17mm,>={Stealth[length=2mm]}]
\draw[FigureGrid,line width=.4pt] (-1.5,-1.9) grid[step=0.5] (1.5,1.9);
\draw[->,FigureStroke,line width=.6pt] (-1.55,0) -- (1.65,0) node[right,text=FigureInk] {$a$};
\draw[->,FigureStroke,line width=.6pt] (0,-1.95) -- (0,2.05) node[above,text=FigureInk] {$b$};
\draw[FigureGreen,line width=1.1pt,dashed] (0,0) ellipse [x radius=1,y radius=1.7321];
\draw[FigureYellow,line width=1.1pt,dash dot] (0,0) ellipse [x radius=0.7071,y radius=1.2247];
\draw[FigureBlue,line width=1.1pt] (0,0) circle [radius=1];
\fill[FigureInk] (1,1) circle (1.2pt) node[above right=-0.3mm,font=\figurefont\scriptsize] {$1+x$};
\foreach \v/\l in {1/1,-1/-1} {\node[font=\figurefont\scriptsize,text=FigureInk,anchor=north east] at (\v,0) {$\l$};}
\node[font=\figurefont\scriptsize,text=FigureInk,anchor=north west] at (0.03,1.70) {$\sqrt3$};
% legend
\begin{scope}[shift={(1.85,1.5)}]
\draw[FigureBlue,line width=1.1pt] (0,0) -- (0.3,0) node[right,text=FigureInk,font=\figurefont\scriptsize] {$\|f\|_{H_K}=1$};
\draw[FigureGreen,line width=1.1pt,dashed] (0,-0.3) -- (0.3,-0.3) node[right,text=FigureInk,font=\figurefont\scriptsize] {$\|f\|_{L^2(dz/2)}=1$};
\draw[FigureYellow,line width=1.1pt,dash dot] (0,-0.6) -- (0.3,-0.6) node[right,text=FigureInk,font=\figurefont\scriptsize] {$\|f\|_{L^2(dz)}=1$};
\end{scope}
\end{tikzpicture}
